\chapter{Discussion and Conclusion}
\label{chap:discussion-conclusion}

\section{Answers to the Research Questions}

\subsection{Answer to RQ1}

\noindent\textbf{RQ1:} \emph{What evidence is required to conclude that a graph learning system has learned a path-aware dependency of a specified length?}
The required evidence must show that the fitted prediction changes with the relevant path information while plausible alternatives are held fixed.
Accuracy alone cannot show this because a system may fit the labels through a shortcut function that does not use the intended dependency.

Four claims must remain separate when interpreting this evidence.
Reachability means that the relevant source information lies inside the target node's receptive field.
Expressibility means that some parameter setting in the model class can realise the target function.
Learnability concerns whether a training procedure can find suitable parameters from finite data and resources.
Demonstrated use means that the fitted prediction actually depends on the path information required by the target function.
Reachability is necessary for standard local message passing, and expressibility removes an impossible target as an explanation of failure.
Neither property establishes learnability or demonstrated use.

The length in the claim must also be defined precisely.
Graph distance is the length of a shortest directed path between two nodes, while witness-path length is the number of edges on the particular path that demonstrates the task dependence.
Dependency length belongs to the target function and states the length over which relevant information must affect the output.
The GLoRa generator parameter $d$ controls a range of realised chain lengths.
Property~(P1) also uses $d$ to index its dependency guarantee, but the parameter is not the exact witness-path or realised chain length of every example.
These quantities may have the same numerical value in a particular case, but none should be substituted for another without establishing that relation.

Paired examples or interventions provide the decisive behavioural evidence.
They must preserve reachability and plausible label-correlated alternatives while changing the relevant source-to-target path information.
The correct output and the fitted prediction should then change in the way required by the path-aware target function.
GLoRa provides this evidence through its path-aware definition and its guarantee over fitting functions under stated sampling conditions \cite{zhou2025glora}.
The conclusion applies only to the fitted function, the specified dependency, and the evaluation conditions.
It does not establish use of every long-range relation or identify why a system fails when the required response is absent.

\subsection{Answer to RQ2}

\noindent\textbf{RQ2:} \emph{How can a benchmark isolate a path-aware dependency of a specified length while ruling out shortcut functions, retaining target expressibility, and comparing systems fairly?}
A benchmark can do so by treating each threat to interpretation with a separate control and by stating the conclusion that each control permits.
The benchmark must define the target output and path-aware dependency, relate the controlled length to realised witness paths, obstruct named shortcut functions, establish target expressibility, and control both class generation and the comparison protocol.

GLoRa implements these requirements for one family of inductive binary node-classification tasks \cite{zhou2025glora}.
Property~(P1) addresses shortcut functions through a sample-qualified guarantee over every fitting node-classification function.
For every generator parameter $d$, precision $\delta>0$, and probability $P\in(0,1)$, there is a sample size $K$.
For this $K$, with probability at least $P$, every function fitting a balanced sample of $K$ generated examples relies on a dependency of length $d$ with precision $\delta$.
The guarantee therefore depends on the sample size, precision, and probability.
Paper~I does not give numerical values of $\delta$ and $P$ that connect the experimental sample size to a particular instance of this guarantee.

Property~(P2) establishes that GCNs can express a fitting function.
Paper~I argues that the other evaluated GNN-based approaches are at least as expressive as GCNs for this task \cite{barcelo2020logical,zhou2025glora}.
This is an expressibility result, not a learnability result, because it proves that suitable parameters exist rather than that optimisation will find them.

Property~(P3) supplies class-distribution fairness by generating positive and negative examples through the same probabilistic algorithm with the requested label selecting the class.
This condition is distinct from fairness across systems.
System-protocol fairness comes from comparing systems on the same task, split, validation procedure, metric, and reporting protocol.

Together, these controls make successful fitting evidence for the specified path-aware dependency under the conditions of Property~(P1), rather than evidence only for prediction on graphs that contain long paths.
Failure can be interpreted as a learning failure only because Property~(P2) establishes that the target is expressible.
The shared empirical protocol makes comparisons between systems interpretable.
These conclusions apply only to GLoRa's dependency family, theorem conditions, generated distributions, evaluated model classes, and experimental protocol.
It does not show that all possible shortcut functions have been excluded from every long-range benchmark or that an expressible target will be learned.

\subsection{Answer to RQ3}

\noindent\textbf{RQ3:} \emph{Under controlled evaluation, how does performance change with dependency length, and to what extent do the diagnostics support common failure explanations?}
Under Paper~I's protocol, accuracy declines as the generator parameter $d$ and the associated realised path lengths increase: most evaluated systems show a performance drop at some tested value of $d$ below $9$, and the stronger systems perform poorly at tested values above $11$ \cite{zhou2025glora}.
These values describe the tested systems and protocol rather than universal dependency-length limits.
The accuracy curve is a symptom of failure to fit the controlled task, not a diagnosis of its mechanism.

Paper~I tests selected traces associated with over-smoothing, over-squashing, and vanishing gradients \cite{li2018deeper,alon2021bottleneck,glorot2010difficulty,zhou2025glora}.
For over-smoothing, it measures only a one-dimensional final target-node representation across examples.
The result shows no complete collapse of that scalar, but it does not establish the absence of all forms of over-smoothing.
For over-squashing, the bounded number of directed source-to-target paths excludes exponential directed-path proliferation in $d$ as a complete explanation.
It does not exclude sensitivity decay, finite-width limits, task-relative mixing, optimisation effects, or all forms of over-squashing.
For vanishing gradients, the inspected weight gradients remain non-zero, but this does not prove useful task-aligned gradient transport.

The diagnostics therefore test only specific traces associated with these three explanations.
They do not identify the cause of the observed errors, and they do not support a universal conclusion about any of the three mechanisms.

\section{Cumulative Contribution}

Paper~I supplies the thesis's empirical contribution.
It introduces GLoRa, states Properties~(P1)--(P3), evaluates the selected graph learning systems, and reports the three diagnostics \cite{zhou2025glora}.
Its contribution is a controlled benchmark and the evidence produced under its formal and empirical conditions.

The kappa contributes the cumulative argument that makes this evidence interpretable.
It separates structural reachability, expressibility, learnability, and demonstrated use.
It defines the evidence needed for a path-aware dependency, compares benchmark designs by the claims their scores can support, and links each proposed failure mechanism to a specific diagnostic trace.
It then places GLoRa's guarantees, controls, performance evidence, and diagnostic limits within that framework.

The cumulative contribution is therefore an evidence standard for claims about one controlled family of long-range graph dependencies.
The standard connects a stated dependency to formal guarantees, empirical controls and comparisons, and bounded diagnostic interpretation.
Because GLoRa is the sole included paper, this contribution is a synthesis around one empirical contribution rather than a portfolio of papers.
It does not provide a general theory of long-range graph learning.

\section{Significance for Evaluation Practice}

The thesis supports a change from reporting long-range performance as one aggregate score to reporting what the score demonstrates.
An evaluation should first define the dependency and its length, then construct evidence that rules out the relevant named shortcut functions.
It should establish that the target function is expressible by each model class for which failure will be interpreted as a learning failure.

Fairness should be reported at two levels.
Class-distribution fairness concerns whether positive and negative examples differ only through the target condition under the stated generator.
System-protocol fairness concerns whether systems are compared under common data, splits, metrics, validation, and reporting conditions.
Combining these two meanings under one fairness claim hides different threats to interpretation.

Performance should be plotted against controlled dependency length, with realised path lengths distinguished from any generator parameter.
When a curve reveals a failing regime, the next step is to test the mechanism-specific trace rather than assign a familiar explanation from accuracy alone.
This practice makes clear which dependency was tested, where the evidence for demonstrated use ends, and what remains unknown about the cause.

Controlled synthetic evidence and real-world evidence remain complementary.
The controlled task can isolate a known dependency, while a real-world task can test practical performance under domain conditions.
Results from one setting should not be used as a substitute for the claims supported by the other.

\section{Limitations}

The first limitation is synthetic control.
GLoRa's guarantees depend on a controlled data-generating process, which does not capture the full range of noise, relation types, missing-edge patterns, repeated motifs, degree distributions, or label-generating processes found in deployed graph domains.
Success or failure on GLoRa therefore does not determine performance on every real-world graph task.

The second limitation is the dependency family.
GLoRa studies one source-to-target path-aware relation in inductive binary node classification.
Tasks based on several distant sources, interacting paths, regional counts, graph cuts, spectral structure, or graph-level relations may require different definitions and controls.

The third limitation is feature simplicity.
The deliberately simple node features and path condition make the task and its shortcut analysis interpretable.
They do not establish how a system behaves when relevant information is distributed across richer node features, edge types, temporal signals, or several path conditions.

The fourth limitation is empirical scope.
Paper~I evaluates a broad but finite set of architectures, training regimes, hyperparameter choices, and data scales under one protocol.
Its results do not rule out better performance from an untested system or protocol, so the empirical conclusions apply only to the reported comparisons.

The fifth limitation is causal incompleteness.
The diagnostic study tests selected traces but does not identify the mechanism that produces the observed errors.
Failure to observe those traces within their measurement bounds does not exclude other forms of the same mechanisms or a different explanation.
Possible explanations involving optimisation, inductive bias, representation loss, or interactions between finite data and finite width remain hypotheses.

\section{Future Work}

To address the synthetic-control limitation, future work should place a dependency with a formal guarantee in real graph structures and use controlled counterexamples to test for shortcut functions.
The dependency-family limitation calls for benchmarks with several sources, interacting paths, ordered path conditions, distant comparisons, or graph-level labels.
Each extension should define its dependency before deriving shortcut, expressibility, and fairness controls for that target.
Richer node features, relation labels, temporal information, or several path conditions would address the feature-simplicity limitation while preserving evidence that a fitted function uses the intended dependency.

Broader evaluation of architectures, training regimes, hyperparameter choices, and data scales under the same controlled protocol would address the empirical-scope limitation.
To address the causal-incompleteness limitation, task-aligned interventions, sensitivity measurements, and gradient-alignment checks should test what information a system preserves and uses.
Both directions should keep their conclusions specific to the systems, settings, and diagnostic traces tested.

Paper~I's relatively strong results for gated-edge systems motivate controlled comparisons of otherwise matched systems with and without gating.
Such comparisons could test whether gating improves the retention of task-relevant information and, in turn, performance.
Until then, the reported association should not be treated as evidence that gating caused the performance difference.

\section{Conclusion}

This thesis establishes an evidence standard for learning a path-aware dependency of specified length and shows how GLoRa supplies the required evidence for one controlled task family.
The evidence combines shortcut control, target expressibility, class-distribution fairness, system-protocol fairness, controlled length, and diagnostics tied to predicted traces.
Within the tested setting, these elements support bounded conclusions about demonstrated use and the inspected failure explanations, but not a universal limit or a complete cause of long-range failure.
The central conclusion is that a long-range score supports a claim about demonstrated use only when the dependency, controls, comparison protocol, and limits of the evidence are stated together.
